\documentclass{article}

\PassOptionsToPackage{numbers}{natbib}
\usepackage[preprint]{neurips_2026}

\usepackage[utf8]{inputenc}
\usepackage[T1]{fontenc}
\usepackage{hyperref}
\usepackage{url}
\usepackage{booktabs}
\usepackage{amsfonts}
\usepackage{amsmath}
\usepackage{amssymb}
\usepackage{nicefrac}
\usepackage{microtype}
\usepackage{xcolor}
\usepackage{graphicx}
\usepackage{enumitem}
\usepackage{listings}
\usepackage{multirow}
\usepackage{array}

\lstset{
  basicstyle=\small\ttfamily,
  breaklines=true,
  frame=single,
  xleftmargin=1em,
}

\definecolor{good}{RGB}{0,128,0}
\definecolor{warn}{RGB}{200,100,0}
\definecolor{note}{RGB}{0,80,160}

\title{Scenario Perturbation Design for Newton--Raphson Surrogate Training:\\
  Current Logic, Physical Rationale, and Future Directions}

\author{Changhun Kim}

\begin{document}
\maketitle

\begin{abstract}
This document describes the scenario perturbation pipeline used to generate
training and test data for a Newton--Raphson (NR) power-flow surrogate model.
We explain in detail what each perturbation parameter controls, why each design
choice is physically motivated, and where the current design still has room for
improvement.  We also draw on the \textsc{gridfm-datakit} library
\cite{gridfm2024} as inspiration for three concrete future directions:
physics-grounded load-scale bounds, admittance perturbation, and improved
topology diversity.
\end{abstract}

%==============================================================================
\section{Background: What the Surrogate Needs to Learn}
%==============================================================================

A Newton--Raphson power-flow surrogate takes as input:
\begin{itemize}[noitemsep]
  \item the grid structure (admittance matrix $Y_{\rm bus}$, bus types),
  \item the power injection schedule $\mathbf{S}^{\rm set} \in \mathbb{C}^N$
        (complex load/generation at every bus), and
  \item an initial voltage guess $\mathbf{u}^{(0)} \in \mathbb{C}^N$,
\end{itemize}
and produces as output the converged voltage solution
$\mathbf{u}^* \in \mathbb{C}^N$.
Training data must therefore contain diverse and physically realistic
\emph{operating points} (different $\mathbf{S}^{\rm set}$) as well as diverse
\emph{starting points} (different $\mathbf{u}^{(0)}$), so that the model
generalises beyond a single base-case operating condition.

%==============================================================================
\section{Overview of the Perturbation Pipeline}
%==============================================================================

Each training sample is generated by the following ordered steps, applied to
a freshly loaded pandapower network:

\begin{enumerate}[noitemsep]
  \item \textbf{Global load scale} --- shift the total system loading level.
  \item \textbf{Per-bus active-power jitter} --- add spatial load diversity.
  \item \textbf{Per-bus reactive-power jitter} --- vary power factors independently.
  \item \textbf{Per-generator active-power jitter} --- vary dispatch setpoints.
  \item \textbf{PV voltage-setpoint jitter} --- vary generator voltage targets.
  \item \textbf{Line outages} --- occasional N-1 topology changes.
  \item \textbf{PPC compilation} --- pandapower compiles the perturbed network
        into a PyPower case (PPC), extracting $Y_{\rm bus}$, bus types, and
        $\mathbf{S}^{\rm set}$.
  \item \textbf{Start voltage construction} --- build $\mathbf{u}^{(0)}$ from a
        DC power-flow solution, then add random angle and magnitude noise.
  \item \textbf{Newton--Raphson solve} --- run NR on the compiled case and record
        the converged solution (or zero it out if NR fails).
\end{enumerate}

Steps 1--6 perturb \emph{what is being solved} (the operating point).
Steps 8--9 control \emph{how the solver starts} (the initial guess diversity).
Step 7 is deterministic given the perturbed network.

%==============================================================================
\section{Perturbation Parameters: Detailed Explanation}
%==============================================================================

\subsection{Global Load Scale (\texttt{load\_scale\_range})}

\paragraph{What it does.}
A single scalar $\alpha \sim \mathcal{U}(l, u)$ is drawn \emph{once per sample}
and applied to every load bus simultaneously:

\begin{equation}
  \tilde{P}^{\rm load}_{i} = \alpha \cdot P^{\rm load}_{i,\,\rm base},
  \qquad
  \tilde{Q}^{\rm load}_{i} = \alpha \cdot Q^{\rm load}_{i,\,\rm base},
  \qquad \forall\, i \in \mathcal{L},
  \label{eq:global_scale}
\end{equation}

where $\mathcal{L}$ is the set of load buses.  When
\texttt{scale\_gen\_with\_load=True} (the default), PV generator
active-power setpoints are scaled by the same factor:

\begin{equation}
  \tilde{P}^{\rm gen}_{g} = \max\!\bigl(\alpha \cdot P^{\rm gen}_{g,\,\rm base},\, 0\bigr),
  \qquad g \in \mathcal{G}_{\rm PV}.
  \label{eq:global_scale_gen}
\end{equation}

\paragraph{Why it is needed.}
Without the global scale, per-bus independent jitter cancels at the aggregate
level.  For a system with $n$ loads and per-bus multiplier
$s_i \sim \mathcal{N}(1, \sigma^2)$, the total load is
$\tilde{P}_{\rm total} = \sum_i P_i s_i$.  By the central limit theorem,
$\tilde{P}_{\rm total} \approx P_{\rm total}$ with a standard deviation of
only $\sigma / \sqrt{n}$, which for $n = 1000$ buses and $\sigma = 0.15$
gives $\approx 0.5\%$ total variation regardless of the per-bus jitter level.
Every sample therefore represents essentially the same operating point.
The global scale breaks this cancellation by moving the entire system to a
different loading regime (light load, nominal, heavy load, near-collapse).

\paragraph{Why gens are scaled together with loads.}
If loads increase but PV generators do not, the slack bus must absorb the
entire excess, which is unrealistic (the slack represents only one generator
in a real system).  Scaling PV generators proportionally ensures that the
global supply--demand balance is maintained before per-generator jitter is
applied; the slack then absorbs only the small residual from per-bus noise,
which is physically realistic.

\subsection{Per-Bus Active-Power Jitter (\texttt{jitter\_load})}

\begin{equation}
  \tilde{P}^{\rm load}_{i} \leftarrow \tilde{P}^{\rm load}_{i}
    \cdot s^P_i, \quad s^P_i \sim \mathcal{N}(1,\, \sigma_P^2),
  \qquad \forall\, i \in \mathcal{L}.
  \label{eq:p_jitter}
\end{equation}

Each load bus receives an \emph{independent} multiplicative draw.
Applied \emph{after} the global scale, this adds spatial diversity:
different buses vary by different amounts around the new global loading level.
The physical motivation is local demand fluctuation --- individual consumers
deviate from the system-wide trend due to stochastic usage patterns.

\subsection{Per-Bus Reactive-Power Jitter (\texttt{jitter\_load\_q})}

\begin{equation}
  \tilde{Q}^{\rm load}_{i} \leftarrow \tilde{Q}^{\rm load}_{i}
    \cdot s^Q_i, \quad s^Q_i \sim \mathcal{N}(1,\, \sigma_Q^2),
  \qquad s^Q_i \perp s^P_i.
  \label{eq:q_jitter}
\end{equation}

Crucially, $s^Q_i$ is drawn \emph{independently of} $s^P_i$.  This allows the
per-bus power factor $\cos\phi_i = P_i / |S_i|$ to vary across samples and
buses.  Using the same multiplier for both P and Q (as in naive implementations)
would freeze the power factor at its base-case value across all samples,
creating a distribution that cannot represent realistic reactive load variation.
In real grids, reactive consumption depends on voltage levels, motor load
fraction, and capacitor switching --- all of which vary somewhat independently
of active power.  We use $\sigma_Q > \sigma_P$ (see Table~\ref{tab:presets})
because reactive power is harder to predict and exhibits higher relative
uncertainty than active power.

\subsection{Per-Generator Active-Power Jitter (\texttt{jitter\_gen})}

\begin{equation}
  \tilde{P}^{\rm gen}_{g} \leftarrow \tilde{P}^{\rm gen}_{g}
    \cdot s^G_g, \quad s^G_g \sim \mathcal{N}(1,\, \sigma_G^2),
  \qquad \forall\, g \in \mathcal{G}_{\rm PV}.
  \label{eq:gen_jitter}
\end{equation}

This models economic dispatch uncertainty: the actual generator output
deviates from the nominal schedule due to forecast errors, ramp constraints,
and real-time re-dispatch.  Note that \eqref{eq:gen_jitter} is applied
\emph{after} the global scaling in \eqref{eq:global_scale_gen}, so the
effective per-generator output is
$P^{\rm gen}_{g,\,\rm base} \cdot \alpha \cdot s^G_g$.

\subsection{PV Voltage-Setpoint Jitter (\texttt{pv\_vset\_range})}

\begin{equation}
  \tilde{V}^{\rm set}_{g} \sim \mathcal{U}(v_{\rm lo},\, v_{\rm hi}),
  \qquad \forall\, g \in \mathcal{G}_{\rm PV},
  \label{eq:vset_jitter}
\end{equation}

where $[v_{\rm lo}, v_{\rm hi}]$ is a \emph{symmetric} interval around
$1.0\,\mathrm{p.u.}$  PV (voltage-controlled) buses hold their terminal
voltage at the generator's setpoint $\tilde{V}^{\rm set}_g$; this jitter
therefore shifts the \emph{solution} --- not just the starting guess.
Different voltage setpoints produce different reactive power flows throughout
the network, resulting in genuinely different converged voltage profiles.

\paragraph{Important: PV vs.\ PQ distinction.}
\eqref{eq:vset_jitter} applies only to PV buses (generators).  PQ buses
(loads, passive nodes) have no voltage setpoint to jitter --- their terminal
voltage is a free variable determined by power-flow physics, not by a
controller.  Confusing these two would be physically incorrect.

\subsection{Line Outages (\texttt{line\_outage\_prob})}

Each in-service line is independently taken out of service with probability
$p_{\rm out}$:

\begin{equation}
  \mathbb{1}[\text{line }k\text{ out}] \sim \mathrm{Bernoulli}(p_{\rm out}),
  \qquad \forall\, k \in \mathcal{E}.
  \label{eq:outage}
\end{equation}

After applying outages, a connectivity check (via \texttt{unsupplied\_buses})
verifies that all buses remain reachable from the slack bus.  If the network
becomes islanded, all outages are restored and the sample proceeds with the
full topology.  This ensures NR always receives a connected network.

Line outages are applied \emph{before} PPC compilation so that the compiled
$Y_{\rm bus}$ and branch-status arrays reflect the actual in-service topology.
This is important: the surrogate must learn to map $(Y_{\rm bus},
\mathbf{S}^{\rm set})$ to $\mathbf{u}^*$ for \emph{different} topologies, not
just for different power injections on a fixed topology.

\subsection{Start Voltage (\texttt{start\_mode}, \texttt{rand\_u\_start},
  \texttt{angle\_jitter\_deg}, \texttt{mag\_jitter\_pq})}

\paragraph{DC compile start.}
With \texttt{start\_mode=dc\_compile}, the initial voltage is constructed from
a DC power-flow solution of the same (perturbed) network:

\begin{equation}
  \theta^{(0)} = \theta^{\rm DC}, \qquad
  |V^{(0)}_{i}| =
  \begin{cases}
    V^{\rm set}_g & \text{if bus }i\text{ is PV or slack} \\
    1.0\,\text{p.u.} & \text{otherwise}
  \end{cases}
  \label{eq:dc_start}
\end{equation}

The DC angles provide a physically informed starting point that reflects the
approximate power-flow solution, which typically requires fewer NR iterations
than a flat start.

\paragraph{Start-point noise.}
When \texttt{rand\_u\_start=True}, random noise is added on top of the DC
solution to increase starting-point diversity:

\begin{align}
  \theta^{(0)}_{i} &\leftarrow \theta^{(0)}_{i} + \delta_i,
  \quad \delta_i \sim \mathcal{U}(-\Delta\theta,\, +\Delta\theta),
  \qquad \forall\, i \notin \{\text{slack}\},
  \label{eq:angle_jitter} \\[4pt]
  |V^{(0)}_{i}| &\leftarrow |V^{(0)}_{i}| \cdot \epsilon_i,
  \quad \epsilon_i \sim \mathcal{U}(1 - \varepsilon_V,\, 1 + \varepsilon_V),
  \qquad \forall\, i \in \mathcal{B}_{\rm PQ}.
  \label{eq:mag_jitter}
\end{align}

Angle noise $\delta_i$ is applied to all non-slack buses (PV and PQ), because
bus angle is a free variable for all non-reference buses.
Magnitude noise $\epsilon_i$ is applied \emph{only to PQ buses}, because:
\begin{itemize}[noitemsep]
  \item PV buses hold their voltage magnitude at $V^{\rm set}_g$ --- the generator
        regulates back to the setpoint anyway, so perturbing the start magnitude
        is meaningless.
  \item The slack bus is the reference node --- its magnitude and angle are fixed
        by definition.
\end{itemize}

\paragraph{Effect on the label ($\mathbf{u}^*$).}
Start-point noise does \emph{not} change the converged solution
$\mathbf{u}^*$.  Its purpose is to diversify the input feature
$\mathbf{u}^{(0)}$, so the surrogate learns to map from
\emph{any} reasonable starting guess to the correct solution.

%==============================================================================
\section{Scenario Presets}
%==============================================================================

Five named presets encode different levels of perturbation intensity.  Their
parameters are summarised in Table~\ref{tab:presets}.

\begin{table}[h]
\centering
\caption{Scenario preset parameters.  $\sigma_P, \sigma_Q, \sigma_G$: Gaussian
  standard deviations for load P, load Q, and generator P jitter.
  $\Delta\theta$: max angle noise [deg] on DC start.
  $\varepsilon_V$: fractional magnitude noise on PQ-bus DC start.
  $p_{\rm out}$: per-line outage probability.}
\label{tab:presets}
\renewcommand{\arraystretch}{1.25}
\begin{tabular}{lccccccccc}
\toprule
Preset & $[l,u]$ & $\sigma_P$ & $\sigma_Q$ & $\sigma_G$
       & $[v_{\rm lo}, v_{\rm hi}]$
       & $\Delta\theta$ & $\varepsilon_V$ & $p_{\rm out}$ \\
\midrule
\texttt{no\_change}
  & ---          & 0     & 0     & 0     & $[1,1]$
  & 0           & 0      & 0 \\
\texttt{easy}
  & $[0.90, 1.10]$ & 0.02  & 0.03  & 0.01  & $[0.998, 1.002]$
  & 0.5°        & 0.002  & 0 \\
\texttt{A}
  & $[0.70, 1.30]$ & 0.05  & 0.08  & 0.03  & $[0.97, 1.03]$
  & 3°          & 0.010  & 0 \\
\texttt{B}
  & $[0.60, 1.40]$ & 0.10  & 0.15  & 0.05  & $[0.95, 1.05]$
  & 5°          & 0.020  & 0.01 \\
\texttt{C}
  & $[0.50, 1.50]$ & 0.15  & 0.15  & 0.08  & $[0.93, 1.07]$
  & 7°          & 0.030  & 0.02 \\
\bottomrule
\end{tabular}
\end{table}

\paragraph{Design rationale for the progression.}
Each step from \texttt{easy} to \texttt{C} widens the perturbation in a
physically motivated way:
\begin{itemize}[noitemsep]
  \item \textbf{Global scale range} widens from $\pm$10\% to $\pm$50\%, covering
    the full range from off-peak (light load) to near-collapse (heavy load).
    The values correspond roughly to: easy = intra-day fluctuations,
    A = day/night swing, B = seasonal variation, C = stress-test.
  \item \textbf{Per-bus jitter} ($\sigma_P$, $\sigma_Q$) is proportionate to
    the global range --- 2\% per-bus noise is consistent with a $\pm$10\%
    global range, while 15\% per-bus noise makes sense in the $\pm$50\%
    stress scenario.
  \item \textbf{$\sigma_Q \geq \sigma_P$}: reactive power is less predictable
    than active power.  We cap $\sigma_Q$ at 0.15 (same as $\sigma_P$ for C)
    because a 22\% Gaussian standard deviation would place one-in-three buses
    at more than 22\% off nominal --- physically unrealistic even for a stress test.
  \item \textbf{PV voltage setpoints} are symmetric around 1.0\,p.u.\ with
    increasing width: $\pm$0.2\% (easy) $\to$ $\pm$7\% (C).  The practical
    operating range for generator setpoints in real grids is approximately
    $[0.95, 1.05]$, which corresponds to the \texttt{B} preset.
  \item \textbf{Line outages} are introduced only from \texttt{B} onward
    (1\% and 2\% per line), because topology changes fundamentally alter the
    problem structure and should be reserved for harder scenarios.
\end{itemize}

%==============================================================================
\section{What Each Perturbation Changes (Summary)}
%==============================================================================

Table~\ref{tab:what_changes} clarifies which aspect of the power-flow problem
each knob affects.

\begin{table}[h]
\centering
\caption{Effect of each perturbation on the power-flow problem.
  ``Operating point'' means the converged solution $\mathbf{u}^*$ changes;
  ``starting point'' means only $\mathbf{u}^{(0)}$ changes.}
\label{tab:what_changes}
\begin{tabular}{lcc}
\toprule
Parameter & Changes $\mathbf{S}^{\rm set}$ / $Y_{\rm bus}$? & Changes $\mathbf{u}^{(0)}$? \\
\midrule
\texttt{load\_scale\_range} ($\alpha$) & \textbf{Yes} --- shifts total loading & Indirectly (DC start adapts) \\
\texttt{jitter\_load} ($\sigma_P$)     & \textbf{Yes} --- spatial P diversity  & Indirectly \\
\texttt{jitter\_load\_q} ($\sigma_Q$)  & \textbf{Yes} --- spatial Q / PF diversity & Indirectly \\
\texttt{jitter\_gen} ($\sigma_G$)      & \textbf{Yes} --- dispatch diversity    & Indirectly \\
\texttt{pv\_vset\_range}               & \textbf{Yes} --- changes PV solution   & Indirectly \\
\texttt{line\_outage\_prob}            & \textbf{Yes} --- changes $Y_{\rm bus}$ & Indirectly \\
\hline
\texttt{rand\_u\_start} / \texttt{angle\_jitter\_deg} & No & \textbf{Yes} \\
\texttt{mag\_jitter\_pq}               & No & \textbf{Yes} (PQ buses only) \\
\bottomrule
\end{tabular}
\end{table}

%==============================================================================
\section{Convergence Diagnostics in the Dataset}
%==============================================================================

Each generated sample records three convergence columns alongside the power-flow
solution:

\begin{center}
\begin{tabular}{lll}
\toprule
Column & Type & Meaning \\
\midrule
\texttt{converged}     & int8    & 1 = NR converged, 0 = did not converge \\
\texttt{final\_misinf} & float64 & $\|\Delta \mathbf{S}\|_\infty$ at last iteration \\
\texttt{nr\_iterations}& int32   & number of NR iterations executed \\
\bottomrule
\end{tabular}
\end{center}

When NR does not converge, the voltage output $\mathbf{u}^{\rm newton}$ is
set to zero, and \texttt{converged=0} flags this row.  Training code should
filter to \texttt{converged=1} rows unless explicitly studying convergence
behaviour.  The \texttt{--drop\_nonconverged} flag can also be used at
generation time to omit failed rows entirely, at the cost of a smaller dataset.

%==============================================================================
\section{Limitations and Future Directions (Inspired by gridfm-datakit)}
%==============================================================================

The \textsc{gridfm-datakit} library \cite{gridfm2024} introduces several
ideas that go beyond the current design.  We discuss three concrete future
directions in order of estimated impact.

\subsection{Physics-Grounded Load-Scale Upper Bound}

\paragraph{Current limitation.}
The global scale bounds $[l, u]$ in Table~\ref{tab:presets} are fixed
\emph{identical} values for all 27 test grids.  However, each grid has a
different stability margin: a small radial distribution network
(\texttt{case33bw}) may become infeasible at $\alpha = 1.15$, while a
heavily meshed transmission network (\texttt{case118}) might not reach its
voltage-collapse boundary until $\alpha = 1.6$ or higher.  Using a fixed
upper bound therefore wastes samples on infeasible cases for weak grids,
and misses the near-collapse region for strong grids.

\paragraph{gridfm-datakit approach.}
Before generating any sample, \textsc{gridfm-datakit} runs an incremental OPF
scan to find the largest $u$ at which a feasible AC-OPF solution still exists:

\begin{equation}
  u^* = \max\bigl\{ u \,:\, \mathrm{OPF}(\alpha \cdot P^{\rm base}) \text{ converges}\bigr\},
  \quad \alpha \in \{1.0, 1.025, 1.05, \ldots, 2.0\}.
  \label{eq:opf_scan}
\end{equation}

The lower bound is then set as $l = (1 - r) \cdot u^*$ for a chosen coverage
fraction $r$ (e.g., $r = 0.4$).  This automatically adapts $[l, u^*]$ to the
physical capacity of each grid, maximising coverage of the feasible operating
manifold without wasting samples in the infeasible region.

\paragraph{Proposed implementation.}
A one-time pre-computation script (per grid, run once before data generation)
would call \texttt{pp.runopp()} at increasing $\alpha$ steps, record $u^*$,
and write a \texttt{load\_scale\_bounds.json} file.  The sbatch scripts would
then pass \texttt{--load\_scale\_lo} and \texttt{--load\_scale\_hi}
from this file rather than using the fixed preset values.

\subsection{Admittance Perturbation}

\paragraph{Current limitation.}
The Y-bus matrix is \emph{fixed} for all samples of a given grid.  The
surrogate therefore learns a mapping conditioned on a single, known admittance
matrix, which limits its ability to generalise to grids with different line
parameters (e.g., due to measurement uncertainty, temperature effects, or
future network reinforcement).

\paragraph{gridfm-datakit approach.}
\textsc{gridfm-datakit} perturbs the resistance $R_k$ and reactance $X_k$ of
each branch $k$ independently:

\begin{equation}
  \tilde{R}_k = R_k \cdot \xi^R_k,\quad
  \tilde{X}_k = X_k \cdot \xi^X_k,\quad
  \xi^R_k,\, \xi^X_k \sim \mathcal{U}(1-\sigma_Y,\, 1+\sigma_Y),
  \label{eq:admittance_pert}
\end{equation}

with typical $\sigma_Y \approx 0.1$--$0.2$.  After perturbing $R$ and $X$,
the Y-bus is recomputed from the modified branch data.  Each sample thus has
a \emph{different} $Y_{\rm bus}$, forcing the surrogate to learn how the
power-flow solution depends on the admittance structure, not just on the power
injections.

\paragraph{Physical motivation.}
Line resistance varies with conductor temperature by approximately
$\pm 10\%$ between winter and summer operating conditions.  Reactance also
varies with transformer tap uncertainty ($\pm 5\%$) and line geometry
tolerances.  Admittance perturbation captures this real-world variability and
additionally serves as a data-augmentation technique that prevents the model
from over-fitting to a single nominal Y-bus.

\paragraph{Implementation note.}
In the existing code, admittance perturbation is straightforward to add to
\texttt{case\_generation\_pandapower}: after PPC compilation, scale
\texttt{branch\_ppc[:, BR\_R]} and \texttt{branch\_ppc[:, BR\_X]} by
independent uniform draws before rebuilding $Y_{\rm bus}$ via
\texttt{build\_Y\_stamped\_from\_ppc}.  A new parameter
\texttt{jitter\_admittance} (scalar $\sigma_Y$) would gate this behaviour.

\subsection{Controlled N-k Topology Sampling}
\label{sec:contingency}

\paragraph{Current limitation.}
The legacy line outage applies independent Bernoulli draws per line.  The
number of simultaneous outages is then Binomially distributed --- mostly zero or
one for small $p_{\rm out}$, but with no explicit control over the contingency
order, and generator outages are not modelled at all.

\paragraph{How much contingency? A survey of the literature.}
Two notions of ``realistic N-$k$'' must not be conflated.  From a
\emph{grid-code} standpoint the mandated criterion is N-1 (the loss of any
single element), defined by NERC TPL-001 \cite{nerc_tpl} in North America and
the ENTSO-E System Operation Guideline in Europe; N-2 is required only for
specific correlated pairs (double-circuit towers, common bus sections,
N-1-1).  This order is essentially \emph{independent of system size} $N$.
From a \emph{probabilistic snapshot} standpoint, however, the number of
elements out at a random instant is $k \sim \mathrm{Poisson}(L\,q)$, where $L$
is the branch count and $q$ the per-branch forced unavailability
($q \approx 3\times10^{-4}$ to $10^{-3}$ from NERC TADS statistics) --- here $k$
\emph{does} scale with $N$.  The risk of high-$k$ contingencies falls sharply
with $k$ and is dominated by a small high-risk subset \cite{chen2005nk}, and
real large-$k$ events are correlated (weather, cascading) rather than
independent draws.

Table~\ref{tab:cont_ratio} surveys the topology mix used by recent ML-for-power-flow
datasets.  Notably most works cap at N-2 and the majority of training data is
N-0 or N-1.

\begin{table}[h]
\centering
\caption{Topology-perturbation mix in recent ML power-flow datasets.}
\label{tab:cont_ratio}
\renewcommand{\arraystretch}{1.2}
\begin{tabular}{lcccc}
\toprule
Dataset & N-0 & N-1 & N-2 & $k_{\max}$ \\
\midrule
\textsc{gridfm-datakit} (default) \cite{gridfm2024} & 0\% & 100\% & 0\% & 1 \\
CANOS (DeepMind) \cite{canos2024}                   & 50\% & 50\% & 0\% & 1 \\
OPFData (DeepMind) \cite{opfdata2024}               & \multicolumn{2}{c}{separate sets} & --- & 1 \\
PF$\Delta$ (MIT) \cite{pfdelta2025}                 & 54.5\% & 27.3\% & 18.2\% & 2 \\
\bottomrule
\end{tabular}
\end{table}

\paragraph{Implemented scheme (\texttt{perturbation\_options.py}).}
We replace the per-line Bernoulli with a selectable
\texttt{sample\_contingency} routine offering three modes:
\begin{itemize}[noitemsep]
  \item \texttt{"bernoulli"} --- legacy behaviour (kept for reproducibility),
        capped at $k_{\max}$.
  \item \texttt{"ratio"} --- PF$\Delta$ style: draw the order $k$ from a fixed
        categorical $\mathbf{w} = (w_0, w_1, w_2)$, default $(0.55, 0.27, 0.18)$.
  \item \texttt{"poisson"} --- size-scaled realism:
        $k \sim \mathrm{Poisson}(L\,q)$, capped at $k_{\max}$, so small grids
        stay near N-0/N-1 and large grids occasionally reach higher orders.
\end{itemize}
In all modes exactly $k$ components are drawn uniformly from the chosen element
set (\texttt{["line"]} or \texttt{["line","gen"]}), connectivity is enforced
(re-draw on islanding, fall back to N-0 after \texttt{max\_tries}), and the
reference/slack machine is never dropped.  Generator outages create a supply
deficit the slack must cover --- physically hard cases that stress NR
convergence.  For the NR surrogate we recommend the PF$\Delta$ ratio
$(0.55,0.27,0.18)$ as the default, with $k_{\max}$ scaled per grid (small grids
$k_{\max}{=}1$, large PEGASE/RTE grids $k_{\max}{=}2$--$3$).

\subsection{Real Temporal Load Profiles (Lower Priority)}

\paragraph{Current limitation.}
The global scale $\alpha \sim \mathcal{U}(l, u)$ assigns equal probability to
all loading levels.  In real grids, load follows a diurnal and seasonal
pattern: overnight load at 50\% is rarer than evening peak at 95\%.
A surrogate trained on uniform sampling may over-represent extreme loading
scenarios relative to their true frequency.

\paragraph{gridfm-datakit approach.}
The global reference is derived from real ERCOT hourly load data, normalised
and scaled to $[l, u^*]$.  Scenarios are indexed by time-step, so the global
scale follows the actual statistical distribution of system load.

\paragraph{Assessment.}
For a surrogate intended to predict NR solutions across the \emph{full feasible
range}, uniform sampling is arguably preferable to a real load profile: the
goal is manifold coverage, not operating-point frequency matching.  A real
profile would concentrate samples at nominal operating conditions and under-represent
the near-collapse region where the surrogate is hardest to train.  We therefore
regard this direction as lower priority than items 1--3 above.

%==============================================================================
\section{Two Dataset Families and the \texttt{share\_grid} Constraint}
%==============================================================================

A subtle but critical implementation constraint governs which perturbations may
coexist in a single dataset.  The training-time loader has a \texttt{share\_grid}
fast path that builds \emph{all} grid-only tensors --- including the dense
$Y_{\rm bus}$ --- once from the first row and reuses them for every sample.  Its
correctness assumption is that \emph{every row describes the same grid}.  This
is what makes large-grid training tractable: without it, every
\texttt{\_\_getitem\_\_} rebuilds a dense complex $Y_{\rm bus}$ (e.g.\ $722^2$
for LVN, $9241^2$ for PEGASE), which dominates I/O.

Perturbations therefore split into two families:

\begin{table}[h]
\centering
\caption{Perturbation families w.r.t.\ the \texttt{share\_grid} optimization.}
\label{tab:families}
\renewcommand{\arraystretch}{1.25}
\begin{tabular}{p{5.5cm}cc}
\toprule
Perturbation & Changes $Y_{\rm bus}$? & \texttt{share\_grid} \\
\midrule
Global load scale, P/Q load jitter, gen jitter & No  & \textbf{True} (safe, fast) \\
PV voltage setpoint, start-point noise         & No  & \textbf{True} (safe, fast) \\
Generator outage (injection / bus-type change) & No  & \textbf{True} (safe, fast) \\
\textbf{Line/transformer outage}               & \textbf{Yes} & \textbf{False} (required) \\
\textbf{Admittance $R$/$X$ jitter}             & \textbf{Yes} & \textbf{False} (required) \\
\bottomrule
\end{tabular}
\end{table}

\paragraph{Why this is a correctness issue, not just speed.}
If a dataset varies $Y_{\rm bus}$ per row but is loaded with
\texttt{share\_grid=True}, every sample silently inherits \emph{row-0's}
$Y_{\rm bus}$ --- the labels no longer match the inputs.  To prevent this, the
\texttt{PerturbationConfig.ybus\_varies} property returns \texttt{True} whenever
a line outage or admittance jitter is active, and the generation script records
this flag in the dataset so the loader selects \texttt{share\_grid} automatically.

\paragraph{Practical consequence for experiment design.}
The two families map onto two experiment tracks that use different grid subsets:
\begin{itemize}[noitemsep]
  \item \textbf{Family 1 (fixed $Y_{\rm bus}$)} --- load/operating-point
        generalization. \texttt{share\_grid=True} makes this affordable on the
        \emph{largest} grids (LVN, PEGASE, RTE). This is the scalability story.
  \item \textbf{Family 2 (varying $Y_{\rm bus}$)} --- topology/parameter
        robustness. \texttt{share\_grid=False}; per-row $Y_{\rm bus}$ is only
        affordable on small-to-medium grids (\texttt{case14}--\texttt{case300}).
        This is the structural-robustness story, and it is exactly where the
        edge-aware attention / known-operator design is differentiated, because
        the per-edge physics input changes with the topology.
\end{itemize}
The two stories therefore do not compete for the same compute: contingency and
admittance experiments live on small/medium grids, while the giant grids carry
the fixed-$Y_{\rm bus}$ scalability experiments.

%==============================================================================
\section{Summary of Current vs.\ Future Design}
%==============================================================================

\begin{table}[h]
\centering
\caption{Comparison of current perturbation design and gridfm-datakit-inspired
  future directions.}
\label{tab:comparison}
\renewcommand{\arraystretch}{1.3}
\begin{tabular}{p{3.2cm}p{4.2cm}p{4.2cm}}
\toprule
\textbf{Aspect} & \textbf{Current} & \textbf{Future direction} \\
\midrule
Global load range
  & Fixed $[l,u]$ identical for all grids
  & Per-grid $u^*$ from OPF feasibility scan \\
\hline
Load P/Q coupling
  & Independent Gaussian noise; $\sigma_Q > \sigma_P$
  & Same; already physically correct \\
\hline
$Y_{\rm bus}$ variation
  & Fixed; never changes within a grid
  & Per-branch $R$/$X$ jitter; Y-bus recomputed \\
\hline
Topology perturbation
  & Bernoulli per-line, lines only
  & Controlled N-$k$, lines + generators \\
\hline
Generator dispatch
  & Proportional scaling + Gaussian jitter
  & OPF-derived dispatch (long term) \\
\hline
Load profile
  & Uniform $\mathcal{U}(l,u)$
  & Real ERCOT temporal profile (low priority) \\
\hline
Convergence tracking
  & \texttt{converged}, \texttt{final\_misinf}, \texttt{nr\_iterations} stored
  & Already implemented \\
\bottomrule
\end{tabular}
\end{table}

\vspace{0.5em}
\noindent
The three highest-priority improvements are, in order:
\begin{enumerate}[noitemsep]
  \item \textbf{Per-grid OPF feasibility scan} for the upper load-scale bound
        --- prevents wasted samples and ensures near-collapse coverage.
  \item \textbf{Admittance perturbation} ($R$/$X$ jitter) --- makes $Y_{\rm bus}$
        vary across samples, teaching the model grid-parameter sensitivity.
  \item \textbf{Controlled N-$k$ topology sampling} including generator
        outages --- creates harder, more physically realistic contingency scenarios.
\end{enumerate}

\vspace{1em}
\begin{thebibliography}{9}
\bibitem{gridfm2024}
M.~Piloto et al.,
\emph{gridfm-datakit: A Python Library for Scalable and Realistic Power Flow
and Optimal Power Flow Data Generation},
arXiv:2512.14658, 2024.

\bibitem{canos2024}
T.~Falconer et al.,
\emph{CANOS: A Fast and Scalable Neural AC-OPF Solver Robust to N-1
Perturbations}, arXiv:2403.17660, 2024.

\bibitem{opfdata2024}
S.~Lovett et al.,
\emph{OPFData: Large-scale datasets for AC optimal power flow with topological
perturbations}, arXiv:2406.07234, 2024.

\bibitem{pfdelta2025}
\emph{PF$\Delta$: A Benchmark Dataset for Power Flow under Load, Generation, and
Topology Variations}, arXiv:2510.22048, 2025.

\bibitem{chen2005nk}
Q.~Chen and J.~D.~McCalley,
\emph{Identifying High Risk N-k Contingencies for Online Security Assessment},
IEEE Transactions on Power Systems, vol.~20, no.~2, 2005.

\bibitem{nerc_tpl}
North American Electric Reliability Corporation,
\emph{TPL-001-4/5: Transmission System Planning Performance Requirements}.
\end{thebibliography}

\end{document}
